%! Introducing Statistics: Do the Data We Collected Tell the Truth? — Unit 3, Lesson 3.1 — Lesson Plan
% GRADUAL-RELEASE plan (warm-up / guided notes and practice / debrief / close and start the
% homework) at the COMPACT budget: the whole lesson is finished in one 55-minute period.
% THERE IS NO GROUP ACTIVITY IN THIS COURSE and NO INDEPENDENT PRACTICE SET in the notes — the
% homework is the individual practice, started in class. Teacher-facing. Every box is `skillbox`
% (breakable) with a \boxguard ahead of it; no `fixedskillbox`, no tier boxes, no exit ticket.
% The first lesson of Unit 3: it opens the unit's question — can we trust how the data were
% collected? — and hands the vocabulary of sampling (3.2–3.4) a reason to exist.
\documentclass[10pt]{article}
\usepackage{apstats-article}
\usepackage{apstats-boxes}

\newcommand{\UnitNumberName}{Unit 3: Collecting Data \quad}
\newcommand{\LessonNumberName}{Lesson 3.1: Introducing Statistics --- Do the Data We Collected Tell the Truth?}

\begin{document}
\begin{center}
    {\Large\bfseries \CourseName \\
    {\small\bfseries \UnitNumberName \LessonNumberName}}
\end{center}

\begin{tcolorbox}[colback=sky, colframe=navy, boxrule=0.9pt, arc=2mm,
    left=2mm, right=2mm, top=1.5mm, bottom=1.5mm]
    \textbf{Primary Objective:} Students will ask who is in a set of data and who decided they
    would be in it, classify a data collection method as chance-based or not, explain in
    context which way a non-chance method pushes a result, and reject the claim that a larger
    sample repairs a non-chance method.
    \\[2pt]
    \textbf{CED alignment:} Topic 3.1 \quad$\cdot$\quad Big Idea: VAR (Variation and
    Distribution) \quad$\cdot$\quad Skill 1.A \quad$\cdot$\quad LO VAR-1.E \quad$\cdot$\quad
    EK VAR-1.E.1 (methods for data collection that do not rely on chance result in
    untrustworthy conclusions)
    \\[2pt]
    \textbf{Lesson model:} \emph{Gradual release --- I do, we do, you do --- all of it inside
    the guided notes.} There is no group activity. The notes name the vocabulary as they teach
    it and deliver the instruction row by row in one \emph{Main Ideas / Notes} table --- each
    idea a definition to complete and a grid of short problems, the first worked by you and the
    rest by the class; the last row, Guided Practice, is worked together with students holding
    the pen; \textbf{the homework is the individual practice}, and students start it in class
    while you circulate. The whole-class debrief is \emph{spoken}.
    \\[2pt]
    \textbf{Compact budget:} the notes are three pages, 14 problems. Every display is a given
    table of methods and results, read and compared; nothing is sketched. The sampling
    vocabulary (population, census, SRS, the named biases) is held back for Lessons 3.2--3.4:
    today the only test is \emph{who decided who is in the data --- chance or people?}
\end{tcolorbox}

\renewcommand{\arraystretch}{1.4}
\boxguard
\begin{skillbox}[Learning Targets \& Key Understandings]{goldbox}
\begin{tabularx}{\linewidth}{@{}X|X@{}}
\textbf{Learning targets (student language)} & \textbf{Key understandings (the ``why'')} \\[2pt]
\hline
\vspace{-6pt}
\begin{itemize}[leftmargin=1.1em, nosep, after=\vspace{2pt}]
  \item I can ask the questions that decide whether data can be trusted: who is in the data,
        and who decided they would be in it?
  \item I can tell a chance-based method from a non-chance method of collecting data.
  \item I can explain in context who a non-chance method over-represents and which way it
        pushes the result.
  \item I can explain why a bigger sample does not fix a non-chance method.
\end{itemize}
&
\vspace{-6pt}
\begin{itemize}[leftmargin=1.1em, nosep, after=\vspace{2pt}]
  \item Methods for data collection that do not rely on chance result in untrustworthy
        conclusions (EK VAR-1.E.1). When people --- the researcher or the respondents ---
        decide who is in the data, some groups are over-represented.
  \item Bias is \emph{consistent}: a non-chance method pushes the result the same direction
        every time it is used, so a careful student can say which way.
  \item \textbf{Target misconception:} ``a bigger sample is a more trustworthy sample.'' Size
        does nothing about \emph{who} is in the data. The news site's 2{,}400 volunteers said
        85\% yes; 80 students chosen at random said 58\%. The 80 win.
  \item \textbf{Second misconception:} ``a lopsided result must mean most people agree.'' The
        skate-shop poll's 88\% says who chose to vote, not what the city thinks.
  \item This lesson sets up the unit: Lessons 3.2--3.4 name the sampling methods and the kinds
        of bias; 3.5--3.7 do the same for experiments.
\end{itemize}
\end{tabularx}
\end{skillbox}

\boxguard
\begin{skillbox}[{Vocabulary, Concepts, \& Theorems --- taught directly in the guided notes}]{greenbox}
\begin{minipage}[t]{0.48\linewidth}
    \begin{itemize}
        \item \textbf{Data collection method} --- how the individuals in the data were chosen
              and measured
        \item \textbf{Chance-based method} --- a random process decides who is in the data
              (EK VAR-1.E.1)
        \item \textbf{Non-chance method} --- people decide: the researcher takes whoever is
              convenient, or respondents volunteer
    \end{itemize}
\end{minipage}
\hfill
\begin{minipage}[t]{0.48\linewidth}
    \begin{itemize}
        \item \textbf{Bias} --- a method that consistently overestimates or underestimates the
              truth (named here informally; Lesson 3.4 names its kinds)
        \item \textbf{The four questions} --- who is in the data? who decided? who is left
              out or over-represented? can it be trusted for everyone?
        \item \textbf{Carried forward} --- explanatory and response variables, and why an
              association is not a cause (Unit 2)
    \end{itemize}
\end{minipage}
\end{skillbox}

% A tabularx never splits, so guard generously ahead of it.
\boxguard[30]
\begin{skillbox}[Lesson at a Glance --- \MeetingLength]{sky}
\renewcommand{\arraystretch}{1.25}
\begin{tabularx}{\linewidth}{@{}l c X X@{}}
\textbf{Phase} & \textbf{Min} & \textbf{Students} & \textbf{Teacher} \\
\hline
Warm-Up & 5 & A sleep-and-scores spiral, the news site's later-start poll, and 2{,}400 or 80
& Hand count on item~3; leave ``2{,}400 --- trustworthy?'' and the tally on the board \\
Guided Notes \& Practice & 34 & Fill the vocabulary box and the notes table: problems 1--10 row
by row, then \emph{The Step Challenge} (11--14) together, pen in their hand & \textbf{I do}
each row's definition lines and first problem (1, 4, 8), \textbf{we do} the rest (24)
$\to$ \textbf{we do} Guided Practice (10) \\
Debrief & 8 & Pens down, papers up --- answer aloud, nothing to fill in & Put 85\% (2{,}400
volunteers) beside 58\% (80 at random); open on the almost-right problem~10 \\
Close \& Assign & 8 & \textbf{Start the homework in class}, alone and silent & Circulate on
Part~A item~3 for your formative read; preview Lesson 3.2 \\
\end{tabularx}
\end{skillbox}

\boxguard[20]
\begin{skillbox}[Warm-Up --- Activate Prior Knowledge (5 min)]{sky}
\begin{minipage}[t]{0.48\linewidth}
    \textbf{The items and what each seeds}
    \begin{itemize}
        \item \textbf{1.} Sleep and test scores: explanatory and response variables, and
              whether an observed association shows cause. \textbf{Spirals} to Unit 2 and
              points ahead to the experiments half of Unit 3.
        \item \textbf{2.} A news site's online poll: 2{,}400 responses, 85\% want a later
              school start. Who answers a poll like this? \textbf{Seeds} the Ask How row ---
              the room's answer becomes \emph{non-chance method}.
        \item \textbf{3.} Trust (A) the 2{,}400 or (B) 80 drawn at random from the roster?
              \textbf{Seeds the crux}: the same pair returns as notes problem~10.
    \end{itemize}
\end{minipage}
\hfill
\begin{minipage}[t]{0.48\linewidth}
    \textbf{Running it}
    \begin{itemize}
        \item \textbf{Debrief item~2 aloud} and write the room's answer to ``who responds?'' on
              the board in their words (``people who care,'' ``people who want it''). Row~1
              names that idea a \emph{non-chance method} ten minutes later.
        \item \textbf{Take a hand count on item~3 and do not settle it.} Expect votes for (A)
              --- ``2{,}400 is a lot.'' Write the tally beside \emph{2{,}400 --- trustworthy?}
              on the board and say problem~10 decides it.
        \item Item~1 is a spiral; take 90 seconds and move on.
    \end{itemize}
\end{minipage}
\end{skillbox}

\boxguard
\begin{skillbox}[Guided Notes \& Practice --- I do, then we do (34 min)]{sky}
\textbf{This is the whole lesson.} There is no group activity, and there is no independent
practice set on this page: \textbf{the homework is the individual practice}, started in class.
The notes are one \emph{Main Ideas / Notes} table --- four rows, 14 short problems.
\textbf{In every row you fill the definition lines and work the first problem of each grid; the
class works the rest, pen in their hand.} The vocabulary box is filled \emph{as you go}.
\begin{multicols}{2}
\textbf{Ask How (6 min).} Fill the three definitions (how / random / people) and read the four
questions aloud --- they are the lesson's procedure, and every later problem is one of them.
Work \textbf{problem~1} yourself (the radio callers chose themselves). Problems 2 and 3 are
theirs; 2 is the quiet one --- the coach is a \emph{researcher} choosing who is handy, which is
still people deciding.

\textbf{Who Is Missing? (8 min).} Fill \emph{consistently} and name \textbf{bias}. Read the
skate-park table top to bottom before anyone writes. Work \textbf{problem~4} (skate-shop fans
are over-represented). \textbf{Problem~6 is the trap}: 88\% looks like a landslide; press for
\emph{why} it is too high --- the fans chose to vote. Problem~7 lands C's 52\% because chance
picked the households, not because it is in the middle.

\textbf{Sample Size (10 min) --- the crux.} Point back at the warm-up on the board. Work
\textbf{problem~8} (the 80 at random). Problem~9 is theirs. \textbf{Problem~10 is the crux}:
take a hand count on ``agree'' \emph{before} anyone explains. Then fill the callout --- a bigger
sample from a non-chance method is still \emph{biased} --- and make the room say it back.

\columnbreak
\textbf{We do (about 10 min): The Step Challenge.} A fresh context, problems 11--14, worked
entirely together. Students hold the pen; you hold the questions: \emph{``Who decided who is
on the leaderboard?''} \quad \emph{``Would a couch-sitter join a step challenge?''}
\textbf{Problem~14 is the crux transferred}: 12{,}000 against 500, and the 500 win. It is the
one to protect when time is short --- a student who answers both 10 and 14 has separated
\emph{how many} from \emph{who}.

\textbf{The pen must actually change hands.} Put problem~11 on them cold, take a hand count on
problem~14 before anyone speaks, and make the room defend the answer. Circulate and demand the
reason:
\begin{itemize}[leftmargin=1.1em, nosep]
  \item \textbf{Problem~12:} ``Who \emph{wants} to be on a leaderboard?''
  \item \textbf{Problem~13:} ``Too high or too low --- and say who pushed it there.''
  \item \textbf{Problem~14:} ``If 50{,}000 joined the challenge, would it get better?''
\end{itemize}
While circulating, pick the papers to put up in the debrief --- one \emph{almost}-right
problem~10 (``disagree, because 80 is enough'') and one that names who volunteered.
\textbf{Your formative read comes twice}: here, and again in the last eight minutes as students
start the homework.
\end{multicols}
\end{skillbox}

\boxguard
\begin{skillbox}[Debrief --- whole class (8 min)]{sky}
Pens down, papers up. \textbf{This debrief is spoken} --- nothing on the page to fill in.
\begin{multicols}{2}
\begin{enumerate}[leftmargin=1.3em, nosep, itemsep=3pt]
  \item \textbf{Open on the almost-right problem~10}: ``disagree, because 80 students is
        plenty.'' Right verdict, wrong reason --- it argues about size, which is the very
        thing that does not matter. Correct it gently and in public: the reason is \emph{who
        decided}.
  \item Put the two pairs side by side on the board: 2{,}400 volunteers, 85\%; 80 at random,
        58\%. Then 12{,}000 on a leaderboard, 11{,}200 steps; 500 at random, 6{,}900. That
        board is the lesson.
  \item Take problem~6 back: did the skate-shop poll tell us about skate parks or about
        skate-shop customers?
\end{enumerate}
\columnbreak
\textbf{The sentence to leave on every desk:} \emph{Before you trust a number, ask who decided
who is in the data. If people decided, more of them does not help.} Make a student say it,
not you.

\textbf{If the debrief runs short}, cut item~3 --- item~1 is the only one that cannot be cut.
\textbf{Do not let the debrief eat the last eight minutes} --- the homework start is your
second formative read.
\end{multicols}
\end{skillbox}

\boxguard
\begin{skillbox}[AP Practice --- extra credit, scored 1 to 5]{goldbox}
Two multiple-choice items --- \textbf{MC~1} picks the chance-based method for a city bus-line
survey (D: 400 random households; the station, website, meeting, and hotline are all people
deciding), and \textbf{MC~2 is the crux in AP clothing}: 50{,}000 streaming-poll responses,
with distractor (A) the target misconception word for word and (D) its cousin --- then one
free-response set on a library lunch survey: \textbf{(a)} did chance decide (no --- the
students chose to be in the library), \textbf{(b)} is 3.2 books too high or too low (too high:
library regulars read more), and \textbf{(c)}, \textbf{the spiral}, mean or median for a
right-skewed distribution (Unit~1).

\textbf{Extra credit, scored 1 to 5, turned in with the homework.} Go over it at
the start of Lesson 3.2. Hold the AP standard out loud: \emph{in context or it does not count}
--- a part (b) that says ``too high because it's biased'' without naming who visits the library
at lunch has not earned the point.
\end{skillbox}

\boxguard
\begin{skillbox}[Homework --- scored 1 to 5, due the first class after two study halls]{goldbox}
Five items in a third context. \textbf{Part~A} (restaurant ratings: 3{,}200 QR-code volunteers
say 41\% excellent, 300 random calls say 63\%): \textbf{A1} which method used chance and who
decided in the other; \textbf{A2} who rates a meal and why 41\% is too low; \textbf{A3} the crux
restated --- ``ten times as many ratings, so report it.'' \textbf{Part~B}: \textbf{B1} an
AP-style multiple choice (the random 600 seniors) with a required justification; \textbf{B2} a
6~a.m.\ gym survey about later hours --- who is over-represented, who is missed, which way it
leans.

\textbf{Start it in the last eight minutes of the period, in class and alone.}
\textbf{Scoring:} 10 points --- A1, A2, A3, B2 at 2 each; B1 1 for the letter and 1 for the
justification. \textbf{A3 is the one to read closely}: a student who sides with the manager has
kept the target misconception, and Lessons 3.3--3.4 build on its absence. The spiralbox
previews Lesson 3.2.

\textbf{Paper homework is the default for this course} --- assign this page unless you have
decided otherwise. \textbf{DeltaMath override (occasional):} on the days you assign DeltaMath
instead, it replaces this page, and this page
stays in the packet as extra practice. Never assign both.
\end{skillbox}

\boxguard
\begin{skillbox}[Watch For (while circulating)]{redbox}
\begin{itemize}[leftmargin=1.2em, nosep]
  \item \textbf{``A bigger sample is more trustworthy''} (notes~10 and 14, HW~A3, AP~MC2 --- the
        target misconception). Probe: \emph{``If 50{,}000 people volunteered, who would they
        be? The same kind of people, only more of them.''}
  \item \textbf{``A lopsided result means most people agree''} (notes~6, AP~(b)). Probe:
        \emph{``Who chose to answer? Does 88\% describe the city or the people who voted?''}
  \item \textbf{Calling any reasonable-sounding method ``random''} (notes~2, HW~B1). The coach
        and the counselors are \emph{people choosing}. Probe: \emph{``Could the result change
        depending on who did the picking? Then it was not chance.''}
  \item \textbf{``Biased'' with no direction} (notes~6 and 13, HW~A2, B2). Probe: \emph{``Too
        high or too low --- and who pushed it there?''}
\end{itemize}
Cold-call prompts: \emph{``Name a method where people decide, but it sounds fair.''} \quad
\emph{``Would you rather have 100 people picked at random or 10{,}000 volunteers? Why?''}
\end{skillbox}

\boxguard
\begin{skillbox}[Close \& Assign (8 min)]{goldbox}
\textbf{Students start the homework here, in class, alone and silent --- this is the individual
practice, and the first minutes of it are your best formative read.} Hand the launch first ---
five items on paper, scored, due the first class after two study halls --- and point out that
the AP Practice page is extra credit, turned in with the homework. Circulate on \textbf{Part~A item~3}. Sort
what you see into three piles as you walk --- disagreed and said the volunteers are still
volunteers (ready); disagreed with no reason, or ``because 300 is enough'' (ready, needs the
language); agreed with the manager (the misconception survived) --- and open Lesson 3.2 by
returning to the 85\%/58\% board if that third pile ran past a third of the room. Then close on
what changed: \emph{``You walked in impressed by 2{,}400. You walk out asking who decided.''}
\textbf{Preview:} Lesson 3.2, \emph{Introduction to Planning a Study} --- the group we want to
know about (the population), the group we actually measure (the sample), and the difference
between watching and experimenting.
\end{skillbox}

% ── Teacher notes — one per component, in packet order (four: there is no activity) ──
\begin{teachernote}[Warm-Up]
Item~3 is the seed and item~2 the quiet setup. Give five minutes total and spend at most 90
seconds on item~1. On item~3 many students will choose (A) because so many people answered; do
not settle it --- take the hand count, write \emph{2{,}400 --- trustworthy?} and the tally on
the board, and say problem~10 decides it. On item~2, write the room's answer to ``who
answers?'' in their words; Row~1 names it a non-chance method ten minutes later.
\end{teachernote}

\begin{teachernote}[Guided Notes \& Practice]
\textbf{Pace it 24 / 10, then 8 for the debrief, then 8 for the homework start.} Four rows, 14
problems. \textbf{You work 1, 4, and 8}; the class works the rest. Ask How must move --- six
minutes --- because the time belongs to the Sample Size row, which carries the crux.
\textbf{Problem~6 is the trap} (88\% reads as consensus); \textbf{problem~10 is the crux}
(2{,}400 volunteers against 80 at random). Take a hand count on 10 before any explanation, and
do not resolve it faster than the room can follow: the callout sentence is the resolution.
\textbf{The Step Challenge is where the pen changes hands}; problem~14 is the crux transferred
and the one to protect. If you only get through 11, 13, and 14, that is the right trade.
\textbf{Your second formative read is the homework start}: Part~A item~3, three piles.
\textbf{Debrief (8 min), spoken}, opening on the almost-right problem~10. \textbf{Do not let the
debrief eat the last eight minutes} --- the homework start is not optional padding.
\end{teachernote}

\begin{teachernote}[AP Practice]
\textbf{This page is extra credit, scored 1 to 5} --- the graded practice is the homework
behind it. \textbf{MC~2 is the one worth reading closely}: (A) is the target misconception
verbatim and (D) is the same idea in a different coat; if more than a handful choose either,
return to the 85\%/58\% board before Lesson 3.3. On the free response, part (b) must give a
direction \emph{and} say who visits the library at lunch; part (c) is a Unit~1 spiral and
should cite the right tail pulling the mean.
\end{teachernote}

\begin{teachernote}[Homework]
\textbf{Scored 1 to 5 (10 points), due the first class after two study halls.} The eight minutes at the
end of the period are a \emph{start}: put students on A1 and A3 first while you can still catch
a wrong start. \textbf{A3 is the item to read first when scoring} --- a student who sides with
the manager has kept the target misconception. On A2, accept any well-argued account of who is
motivated to rate, but the answer must give the direction (too low). On B1, (D) with no
justification earns 1 of 2. On B2, both groups and a direction are needed for full credit.
\textbf{Paper is the default.} On the occasions you assign DeltaMath in its place, tell
students this page is now optional practice. Do not assign
both.
\end{teachernote}
\end{document}
